\section{Decision Tree tự cài đặt và đối chiếu với sklearn}

Ở phần mở rộng này, nhóm tự cài đặt lại các bước cốt lõi của CART bằng NumPy để
kiểm tra mức độ minh bạch của thuật toán, đồng thời đối chiếu với
\texttt{DecisionTreeClassifier} của scikit-learn.

\subsection{Cài đặt từ đầu và protocol so sánh}

\texttt{CustomDecisionTreeClassifier} không gọi lại implementation nội bộ của sklearn.
Model tự cài đặt tìm split nhị phân bằng cách quét tất cả feature và các ngưỡng liền kề
phân biệt, chọn split có Gini impurity nhỏ nhất. Các điều kiện dừng gồm
\texttt{max\_depth}, \texttt{min\_samples\_split} và \texttt{min\_samples\_leaf}; model cũng
trả về class probability và văn bản các nhánh quyết định để phục vụ diễn giải.

Nhóm huấn luyện hai model trên cùng bộ Handwritten Digits, cùng split stratified 80/20
\texttt{random\_state=42}, criterion Gini, \texttt{max\_depth=8} và
\texttt{min\_samples\_leaf=1}. Kết quả được tính trên 360 mẫu test và không dùng test set
để chọn tham số.

\begin{table}[H]
    \centering
    \small
    \caption[So sánh custom tree và sklearn]{So sánh custom Decision Tree và sklearn trên Handwritten Digits.}
    \label{tab:custom-vs-sklearn}
    \setlength{\tabcolsep}{3.5pt}
    \begin{tabularx}{\textwidth}{@{}>{\raggedright\arraybackslash}Xrrrrrrr@{}}
        \toprule
        \textbf{Model} & \textbf{Train acc.} & \textbf{Test acc.} & \textbf{Error} &
        \textbf{Macro-F1} & \textbf{Depth} & \textbf{Leaves} & \textbf{Fit (s)} \\
        \midrule
        Custom tree & 93.95\% & 82.22\% & 17.78\% & 82.10\% & 8 & 82 & 0.170 \\
        sklearn tree & 93.88\% & 80.83\% & 19.17\% & 80.72\% & 8 & 81 & 0.010 \\
        \bottomrule
    \end{tabularx}
\end{table}

Hai model có test prediction trùng nhau 94.72\%. Custom tree có accuracy và Macro-F1
cao hơn trong lần chạy này, nhưng không đủ để kết luận rằng implementation tự cài đặt
tốt hơn sklearn. Các split có cùng mức giảm Gini có thể được xử lý theo thứ tự khác nhau,
dẫn đến topology và prediction khác nhau. sklearn vẫn nhanh hơn rõ ràng nhờ implementation
đã tối ưu; kết quả này phân biệt mục tiêu học tập của custom tree với mục tiêu sản xuất.
Artifact chi tiết được lưu tại \nolinkurl{results/custom_vs_sklearn_decision_tree.json} và bảng
so sánh tại \nolinkurl{results/custom_vs_sklearn_decision_tree__comparison.csv}.

\subsection{Xác nhận pipeline end-to-end}

Để kiểm tra khả năng chạy thật, nhóm chạy pipeline local qua đầy đủ các bước
\textit{load -> validate/preprocess -> split -> fit -> predict -> metrics}. Letter
Recognition được đọc từ split đã làm sạch; Handwritten Digits và Covertype được đọc từ
raw CSV, kiểm tra schema và missing values trước khi chia stratified. Baseline
\texttt{DecisionTreeClassifier} sau đó được fit và đánh giá trên cả ba dataset.

Lần chạy local của nhóm hoàn tất thành công cho 18,668 mẫu Letter Recognition, 1,797 mẫu
Handwritten Digits và 581,012 mẫu Covertype; tất cả các bước đều pass trong
\nolinkurl{results/end_to_end_pipeline_validation.json}. Kết quả test accuracy lần lượt là
86.77\%, 82.50\% và 93.89\%. Đây là bằng chứng cho pipeline CPU của ba dataset trong
project. Notebook benchmark bốn model sử dụng RAPIDS/cuML trên GPU vẫn là artifact phụ
thuộc môi trường Kaggle và không được coi là một lần xác nhận GPU local.
